
\subsubsection*{5.1.3}
\phantomsection
\label{sga3:VIIB:5.1.3}

Identify every profinite \(k\)-algebra \(\Gamma\) with the graded
profinite \(k\)-algebra \((\Gamma_n)_{n\geq0}\) such that
\[
  \Gamma_0=\Gamma,
  \qquad
  \Gamma_n=0\quad\text{if }n>0.
\]
In particular, if \((A_n)_{n\geq0}\) is a graded profinite
\(k\)-algebra, we shall regard \(A_0\) indifferently as a profinite
\(k\)-algebra or as a graded profinite \(k\)-algebra. We shall then
denote by
\[
  \rho:(A_n)\longrightarrow A_0
\]
the morphism of \(\mathbf{Alpg}/k\) such that
\[
  \rho_0=\operatorname{id}_{A_0},
  \qquad
  \rho_n=0\quad\text{if }n>0.
\]
Likewise,
\[
  \tau:A_0\longrightarrow(A_n)
\]
will denote the section of \(\rho\) such that
\[
  \tau_0=\operatorname{id}_{A_0},
  \qquad
  \tau_n=0\quad\text{if }n>0.
\]

Every cogroup structure on \((A_n)_{n\in\mathbb N}\) induces a cogroup
structure on \(A_0\) such that \(\rho\) and \(\tau\) are homomorphisms of
cogroups. Since \(\tau\) is a section of \(\rho\), the cogroup
\((A_n)_{n\in\mathbb N}\) is therefore a ``semidirect sum'' of \(A_0\)
and of a cogroup \((A'_n)_{n\in\mathbb N}\) of \(\mathbf{Alpg}/k\) such
that \(A'_0=k\). In other words, if \((A'_n)\) is the graded subalgebra
of \((A_n)\) such that the sequence
\[
  \begin{tikzcd}[column sep=3.4em]
    (A'_n) \arrow[r,"j=\mathrm{incl.}"]
    & (A_n)
      \arrow[r,shift left=0.7ex,"\mathrm{in}_2"]
      \arrow[r,shift right=0.7ex,"(\rho\widehat\otimes\operatorname{Id})\Delta"']
    & A_0\widehat\otimes_k(A_n)
  \end{tikzcd}
\]
is exact, the morphism
\[
  A_0\widehat\otimes_k(A'_n)\longrightarrow(A_n)
\]
with components \(\tau\) and \(j\) is invertible. \textup{(One denotes
by \(\Delta:(A_n)\to(A_n)\widehat\otimes_k(A_n)\) the morphism which
defines the cogroup structure of \((A_n)\).)}\footnote{Translator's note: \(\Delta\) maps
\((A'_n)\) into \((A'_n)\widehat\otimes_k(A_n)\), but not into \((A'_n)\widehat\otimes_k(A'_n)\) in general,
so the cogroup structure of \((A'_n)\) is not induced by \(\Delta\). The Gille--Polo re-edition (2011),
Expos\'e~VII\textsubscript{B}, makes it explicit; in our formulation: \(\rho\) and \(\tau\) correspond to a
closed immersion of the formal group \(\operatorname{Spf}A_0\) into the formal group of \((A_n)\) and to a
retraction of it; the kernel \(N\) of this retraction has as affine algebra the quotient of \((A_n)\) by the
closed ideal generated by \(\tau(I_0)\), \(I_0\) being the augmentation ideal of \(A_0\); the composite of
\(j\) and of this projection is an isomorphism of graded algebras, through which \((A'_n)\) receives the
cogroup structure of \(N\); and the formal group of \((A_n)\) is the semidirect product of \(N\) by
\(\operatorname{Spf}A_0\). The 1970 text writes \((A_n)_{n>0}\) for \((A_n)_{n\geq0}\) above.}
